% ================================================================
% ORGANIZZATIVI
% ================================================================
\section{Processi Organizzativi}
 
\subsection{Processo Di Gestione}
\label{sec:gestione}
 
\subsubsection{Descrizione}
 
Il processo di gestione regola la pianificazione delle attività, l'assegnazione e
la rotazione dei ruoli, la conduzione delle riunioni e il tracciamento
dell'avanzamento del lavoro. Esso fornisce la struttura organizzativa entro cui
tutti gli altri processi operano.
 
\subsubsection{Ruoli Del Gruppo}
\label{sec:ruoli}
 
Il gruppo prevede i seguenti ruoli, assegnati a rotazione tra i membri. Il
Regolamento del Progetto Didattico impone che ogni componente:
 
\begin{itemize}
  \item ruoti su tutti i ruoli secondo regole documentate;
  \item assuma ogni ruolo per un tempo totale congruo;
  \item assuma \textbf{non più di un ruolo alla volta}.
\end{itemize}
 
Le regole di rotazione devono assicurare assenza di conflitti di interesse,
garantendo verifica indipendente per ogni attività.
 
\bigskip
\begin{tabularx}{\textwidth}{l c X}
  \toprule
  \textbf{Ruolo} & \textbf{Costo orario} & \textbf{Responsabilità} \\
  \midrule
  Responsabile
    & 30\,€/h
    & Coordina l'elaborazione di piani e scadenze; approva il rilascio di
      prodotti parziali o finali (SW, documenti); coordina le attività del
      gruppo; rappresenta \textit{Synergo Unit} verso committente e
      proponente; redige il verbale interno di ogni riunione; redige e presenta il Diario di Borgo settimanale. \\[6pt]
  Amministratore
    & 20\,€/h
    & Assicura l'efficienza di procedure, strumenti e tecnologie a supporto
      del \textit{way of working\textsubscript{G}}; redige i verbali esternie mantiene aggiornate le Norme di
      Progetto, il Piano di Progetto e il Piano di Qualifica; gestisce la creazione delle issue tramite Google Form; verifica,
      prima della chiusura di ogni sprint, che i tempi effettivi nel foglio di
      calcolo siano stati compilati correttamente. \\[6pt]
  Analista\textsubscript{G}
    & 25\,€/h
    & Svolge le attività di analisi dei requisiti. \\[6pt]
  Progettista\textsubscript{G}
    & 25\,€/h
    & Svolge le attività di progettazione (\textit{design}). \\[6pt]
  Programmatore\textsubscript{G}
    & 15\,€/h
    & Svolge le attività di codifica; attivato nella fase PB. \\[6pt]
  Verificatore
    & 15\,€/h
    & Svolge le attività di verifica di documenti e codice prodotti dagli
      altri membri. \\
  \bottomrule
\end{tabularx}
 
\bigskip
Le assegnazioni di ruolo sono tracciate nei verbali interni di ogni sprint e
nei task su GitHub Project\textsubscript{G}.
 
\subsubsection{Pianificazione E Tracciamento Delle Attività}

Le attività sono tracciate tramite \textbf{GitHub Project}.
Ogni task corrisponde a una issue GitHub il cui identificativo numerico funge da
codice univoco. Ogni task riporta obbligatoriamente i seguenti attributi, compilati
tramite il Google Form di creazione task:

\begin{itemize}
  \item \textbf{Tipo}: \textit{doc.\ gestionale}, \textit{doc.\ tecnico}, o
        \textit{altro} (per task non legati a un documento, es.\ aggiornamento sito,
        invio email);
  \item \textbf{Categoria}: \textit{interno}, \textit{esterno}, o \textit{altro};
  \item \textbf{Nome documento}: per i verbali nel formato
        \texttt{verb-int/est-AAAA-MM-GG}; per task non documentali, parole
        identificative senza spazi (es.\ \texttt{campi-google-form},
        \texttt{aggiornamento-sito}). Il nome documento coincide con la parte
        finale del nome del branch;
  \item \textbf{Descrizione}: sintetica; il dettaglio è nel documento di riunione
        condiviso;
  \item \textbf{Priorità}: \textit{Bassa}, \textit{Media}, \textit{Alta},
        \textit{Critica};
  \item \textbf{Assegnazione}: uno o più membri del gruppo;
  \item \textbf{Ruolo}: ruolo dell'assegnatario per quel task;
  \item \textbf{Scadenza};
  \item \textbf{Tempo previsto\textsubscript{G}} per il Redattore (in minuti);
  \item \textbf{Tempo previsto per il Responsabile e l'Amministratore}: se sono i diretti assegnatari, il campo
        è pari a 0;
  \item \textbf{Tempo previsto per il Verificatore};
  \item \textbf{Tempo effettivo} per ciascun ruolo coinvolto (in minuti),
        compilato \textbf{dopo il merge e la chiusura della issue};
  \item \textbf{Stato di avanzamento}: \textit{Backlog}, \textit{In Progress},
        \textit{Done}.
\end{itemize}

La creazione di un task avviene tramite \textbf{Google Form}, che genera
automaticamente una issue su GitHub e registra i dati in un
foglio di calcolo condiviso. Il form è configurato con campi precompilati che
rispettano automaticamente le convenzioni di nomenclatura di task e branch.

I task non relativi a modifiche della repository
vengono classificati come tipo \textit{altro}.

\subsubsection{Calcolo Ore E Costi Sprint}
\label{sec:calcolo-ore}

Il foglio di calcolo condiviso calcola automaticamente, selezionando la data di
inizio sprint, preventivo e consuntivo\textsubscript{G} per lo sprint corrente e le risorse residue\textsubscript{G}.
Le formule di conversione applicate sono:

\begin{align*}
  \text{ore} &= \frac{\text{minuti}}{60} \\
  \text{costo} &= \frac{\text{minuti}}{60} \times \text{costo orario del ruolo}
\end{align*}

Il foglio produce le seguenti tabelle:

\paragraph{Preventivo}
Per ogni membro: ruolo nello sprint (recuperato dalle task), ore preventivate e
costo preventivato, con totali di colonna.

\paragraph{Consuntivo}
Per ogni membro: ruolo, ore preventivate, ore effettive, delta ore, costo
preventivato, costo effettivo, delta costo, con totali di colonna.

\paragraph{Risorse Residue Per Ruolo}
Per ogni ruolo: ore residue, costo orario (fisso per ruolo, cfr.\
Sezione~\ref{sec:ruoli}), costo totale residuo.

\paragraph{Ore Residue Per Membro}
Per ogni membro: ore residue suddivise per ruolo (Responsabile, Amministratore,
Analista, Progettista, Programmatore, Verificatore) e totale complessivo.
 
\subsubsection{Riunioni Interne}
 
\paragraph{Cadenza E Piattaforma}
Il gruppo si riunisce con cadenza fissa, salvo diversa comunicazione del
Responsabile:
 
\begin{itemize}
  \item \textbf{Martedì} alle \textbf{14:30};
  \item \textbf{Venerdì} alle \textbf{14:30}.
\end{itemize}
 
Le riunioni si svolgono in modalità remota sulla piattaforma Discord\textsubscript{G}.
Il Responsabile può convocare riunioni straordinarie tramite il gruppo
WhatsApp\textsubscript{G}.
 
\paragraph{Struttura Della Riunione Del Martedì}
La riunione del martedì segue la struttura seguente, nell'ordine indicato:

\begin{enumerate}
      \item \textbf{Sprint Retrospective\textsubscript{G}}: il team identifica criticità nel processo
            di lavoro e definisce almeno un'azione correttiva concreta da applicare
            allo sprint successivo.
      \item \textbf{Sprint Review\textsubscript{G}}: il team verifica
            l'avanzamento rispetto alle attività pianificate nello sprint appena
            concluso, analizza le eventuali difficoltà e confronta le ore previste
            con quelle effettive per ciascun task.
      \item \textbf{Sprint Planning\textsubscript{G}}: il Responsabile aggiorna il backlog prioritario,
            stima il tempo previsto per i task completabili nello sprint corrente e
            l'Amministratore assegna i task tramite Google Form. Output: sprint backlog.
      \item \textbf{Predisposizione del documento di riunione}: lo Scriba
            registra le decisioni prese nel documento Google condiviso, che costituisce
            la base per la redazione del verbale.
      \item \textbf{Redazione del verbale}: il Responsabile redige il verbale interno
            al termine della riunione; il verbale è verificato e pubblicato entro la
            riunione successiva.
\end{enumerate}

\paragraph{Aggiornamento Asincrono E Riunione Del Venerdì}
Il daily stand-up\textsubscript{G} è gestito in modalità asincrona tramite il gruppo WhatsApp:
ogni membro aggiorna il gruppo su ciò che ha svolto, ciò che prevede di fare e
gli eventuali impedimenti. La riunione del venerdì alle 14:30 funge da punto di
sincronizzazione intermedio per affrontare dubbi o criticità emerse durante la
settimana.
 
\paragraph{Modalità Decisionale}
Il gruppo adotta il \textbf{brainstorming\textsubscript{G}} come tecnica principale per la generazione
e il confronto delle idee. Il processo si articola in tre fasi:
 
\begin{enumerate}
  \item raccolta libera delle proposte da parte di ogni membro, senza valutazione
        immediata;
  \item confronto strutturato delle idee emerse, con analisi di vantaggi e
        svantaggi;
  \item scelta della soluzione ritenuta più adeguata, con registrazione
        della decisione nel verbale.
\end{enumerate}
 
Le decisioni adottate sono registrate con il codice \texttt{VI-y.x}. Le decisioni
non possono essere modificate unilateralmente: eventuali revisioni richiedono
delibera in una riunione successiva.
 
\subsubsection{Gestione Delle Disponibilità}
 
Le disponibilità dei membri e la pianificazione delle riunioni ricorrenti sono
gestite tramite \textbf{Google Calendar\textsubscript{G}}. I calendari condivisi
sono aggiornati dal Responsabile o dall'Amministratore. Le comunicazioni urgenti
e le notifiche rapide avvengono tramite il gruppo \textbf{WhatsApp}.
 
% ---------------------------------------------------------------
\subsection{Processo Di Infrastruttura}
\label{sec:infrastruttura}
 
\subsubsection{Descrizione}
 
Il processo di infrastruttura definisce, mantiene e documenta l'insieme degli
strumenti adottati dal gruppo a supporto di tutti i processi attivi. L'adozione
di un nuovo strumento deve essere deliberata collegialmente e registrata nel verbale
della riunione in cui viene approvata; le presenti norme devono essere aggiornate
di conseguenza.
 
\subsubsection{Strumenti Di Gestione}
 
\begin{itemize}
  \item \textbf{GitHub Project}: sistema principale per la
        pianificazione e il tracciamento dei task, con gestione
        di assegnatari, scadenze, stime e stati di avanzamento.
  \item \textbf{Google Form}: interfaccia guidata per la creazione dei task;
        genera automaticamente la issue corrispondente su GitHub
        e registra i dati nel foglio di calcolo condiviso.
  \item \textbf{Foglio Di Calcolo}: aggregazione e monitoraggio delle misurazioni
        relative ai task; calcola automaticamente preventivo, consuntivo e risorse
        residue per sprint (cfr.\ Sezione~\ref{sec:calcolo-ore}).
  \item \textbf{Google Calendar}: gestione delle disponibilità
        dei membri e pianificazione delle riunioni ricorrenti e straordinarie.
\end{itemize}
 
\subsubsection{Strumenti Di Documentazione}
 
\begin{itemize}
  \item \textbf{\LaTeX{} con Visual Studio Code}: ambiente di redazione documentale
        principale; la compilazione automatica in PDF è gestita tramite l'estensione
        \textit{LaTeX Workshop}.
  \item \textbf{Script Python} (\texttt{python\_glossary\_automation/}):
        automatizza l'applicazione del pedice G ai termini del
        Glossario, riducendo il rischio di omissioni.
  \item \textbf{GitHub}: sistema di controllo versione centralizzato per tutti i
        documenti del progetto; gestisce branch, pull
        request e baseline.
  \item \textbf{GitHub Pages}: piattaforma di pubblicazione del
        sito documentale del gruppo, sviluppato in HTML5 e CSS3.
\end{itemize}
 
\subsubsection{Strumenti Di Comunicazione}
 
\begin{itemize}
  \item \textbf{Email Istituzionale Del Gruppo}: canale ufficiale per le
        comunicazioni formali con committente e proponente. Configurata con
        inoltro automatico verso i membri delegati; le risposte mantengono
        l'indirizzo istituzionale come mittente.
  \item \textbf{Discord}: piattaforma principale per le riunioni
        interne, le discussioni operative e il coordinamento quotidiano tra i
        membri.
  \item \textbf{WhatsApp}: canale secondario per comunicazioni
        rapide, notifiche urgenti e convocazioni straordinarie.
\end{itemize}

\subsection{Processo Di Miglioramento}
\label{sec:miglioramento}
 
\subsubsection{Descrizione}
 
Il processo di miglioramento ha lo scopo di identificare inefficienze
nei processi attivi e adottare azioni correttive concrete, applicando
il ciclo PDCA\textsubscript{G} descritto nella
Sezione~\ref{sec:qualita}.
 
\subsubsection{Attività}
 
\paragraph{Identificazione}
Le criticità emergono da due fonti principali:
 
\begin{itemize}
  \item la \textbf{Sprint Retrospective\textsubscript{G}}: ogni martedì il
        team identifica almeno un'azione correttiva concreta da
        applicare allo sprint\textsubscript{G} successivo, registrata nel
        verbale interno.
  \item il \textbf{monitoraggio delle metriche}: scostamenti significativi
        tra preventivo e consuntivo orario o tra versione attesa e
        versione consegnata segnalano processi da rivedere.
\end{itemize}
 
\paragraph{Attuazione}
L'Amministratore\textsubscript{G} aggiorna le presenti norme in seguito a ogni
delibera di modifica adottata in riunione. Nessun processo può essere
modificato unilateralmente: la revisione richiede delibera esplicita e
registrazione nel verbale interno con codice decisione
\texttt{VI-y.x\textsubscript{G}}.

\subsection{Processo Di Formazione}
\label{sec:formazione}
 
\subsubsection{Descrizione}
 
Il processo di formazione ha lo scopo di garantire che ogni membro del
gruppo acquisisca le competenze necessarie per operare efficacemente
sulle tecnologie e gli strumenti adottati, minimizzando i rischi legati
alla rotazione dei ruoli\textsubscript{G}.
 
\subsubsection{Tecnologie Di Riferimento}
 
Le tecnologie oggetto di formazione per la fase RTB\textsubscript{G} sono:
 
\bigskip
\begin{tabularx}{\textwidth}{l X}
  \toprule
  \textbf{Ambito} & \textbf{Risorsa principale} \\
  \midrule
  Vue.js (frontend\textsubscript{G}) &
    \href{https://vuejs.org/guide/introduction.html}{\underline{Documentazione ufficiale Vue.js}} \\
  CodeMirror (editor) &
    \href{https://codemirror.net/docs/}{\underline{Documentazione ufficiale CodeMirror}} \\
  Python / FastAPI\textsubscript{G} (backend\textsubscript{G}) &
    \href{https://fastapi.tiangolo.com/}{\underline{Documentazione ufficiale FastAPI}} \\
  PostgreSQL\textsubscript{G} (database) &
    \href{https://www.postgresql.org/docs/}{\underline{Documentazione ufficiale PostgreSQL}} \\
  Docker\textsubscript{G} (deployment) &
    \href{https://docs.docker.com/get-started/}{\underline{Docker Get Started}} \\
  \LaTeX{} / LaTeX Workshop &
    \href{https://www.latex-project.org/}{\underline{Documentazione ufficiale \LaTeX{}}} \\
  GitHub\textsubscript{G} &
    \href{https://docs.github.com/}{\underline{GitHub Docs}} \\
  LucidSpark\textsubscript{G} (diagrammi UML\textsubscript{G}) &
    \href{https://lucid.app/}{\underline{Lucidchart / LucidSpark Help}} \\
  \bottomrule
\end{tabularx}
 
\subsubsection{Modalità Di Formazione}
 
\begin{itemize}
  \item \textbf{Auto-formazione}: ogni membro studia autonomamente la
        documentazione ufficiale degli strumenti assegnati prima di
        operare su di essi nel repository\textsubscript{G}.
 
  \item \textbf{Apprendimento tra pari}: in caso di introduzione di una
        nuova tecnologia o di rotazione di ruolo\textsubscript{G}, il membro più
        esperto organizza una breve sessione sincrona su
        Discord\textsubscript{G} per condividere best~practice operative.
 
  \item \textbf{Supporto alla rotazione}: chi assume un ruolo nuovo si
        confronta con chi lo ha ricoperto precedentemente per
        ricevere indicazioni operative non formalmente codificabili.
\end{itemize}
 
L'adozione di un nuovo strumento non presente nella tabella sopra
deve essere deliberata in riunione, registrata nel verbale interno
e aggiornata nelle presenti norme prima dell'utilizzo effettivo
(principio \textit{just-in-time\textsubscript{G}}).